\documentclass[10pt,a4paper,fontset=none]{ctexart}
\usepackage[margin=1.65cm]{geometry}
\usepackage{xcolor}
\usepackage{listings}
\usepackage{longtable}
\usepackage{booktabs}
\usepackage{array}
\usepackage{enumitem}
\usepackage{hyperref}
\usepackage{fancyhdr}
\usepackage{titlesec}
\usepackage{tcolorbox}
\tcbuselibrary{breakable}
\usepackage{fontspec}
\setCJKmainfont{Songti SC}
\setCJKsansfont{Songti SC}
\setCJKmonofont{Songti SC}
\setmonofont{Menlo}[Scale=0.82]

\hypersetup{colorlinks=true, linkcolor=blue!45!black, urlcolor=blue!60!black}
\setlength{\headheight}{14pt}
\pagestyle{fancy}
\fancyhf{}
\lhead{DSA 期末考前学案}
\rhead{模板迁移与作业复盘}
\cfoot{\thepage}
\titleformat{\section}{\Large\bfseries\color{blue!45!black}}{\thesection}{0.7em}{}
\titleformat{\subsection}{\large\bfseries\color{black}}{\thesubsection}{0.6em}{}
\titleformat{\subsubsection}{\normalsize\bfseries\color{blue!35!black}}{\thesubsubsection}{0.5em}{}
\setlist[itemize]{itemsep=0.18em, topsep=0.25em, leftmargin=1.7em}
\setlist[enumerate]{itemsep=0.18em, topsep=0.25em, leftmargin=1.9em}

\definecolor{codebg}{RGB}{248,248,248}
\definecolor{darkgreen}{RGB}{0,100,0}
\lstdefinestyle{py}{
  language=Python,
  backgroundcolor=\color{codebg},
  basicstyle=\ttfamily\scriptsize,
  keywordstyle=\color{blue!70!black}\bfseries,
  commentstyle=\color{darkgreen},
  stringstyle=\color{red!55!black},
  showstringspaces=false,
  frame=single,
  rulecolor=\color{black!20},
  breaklines=true,
  breakatwhitespace=false,
  tabsize=4,
  columns=fullflexible,
  keepspaces=true
}
\lstdefinestyle{plain}{
  backgroundcolor=\color{codebg},
  basicstyle=\ttfamily\scriptsize,
  frame=single,
  rulecolor=\color{black!15},
  breaklines=true,
  columns=fullflexible,
  keepspaces=true
}

\newtcolorbox{focusbox}{colback=blue!3!white,colframe=blue!40!black,title=本节核心,breakable}
\newtcolorbox{warnbox}{colback=red!2!white,colframe=red!45!black,title=高频失误,breakable}
\newtcolorbox{methodbox}{colback=green!2!white,colframe=green!40!black,title=迁移方法,breakable}
\newtcolorbox{checkbox}{colback=yellow!5!white,colframe=yellow!45!black,title=考前检查,breakable}
\newcolumntype{L}[1]{>{\raggedright\arraybackslash}p{#1}}

\title{\bfseries DSA 期末考前学案：从模板到实例的迁移\\\large 结合错题 PDF、tree.md 与 graph.md 的详细复盘}
\author{根据当前文件夹资料整理}
\date{2026 年 6 月}

\begin{document}
\maketitle
\tableofcontents
\newpage

\section*{使用说明}
\addcontentsline{toc}{section}{使用说明}
这份学案不是把 \texttt{tree.md} 和 \texttt{graph.md} 再抄一遍，而是围绕一个真实考试问题展开：模板代码能看懂，但换成题目实例之后，不知道该改哪里。复习时建议按如下顺序使用：
\begin{enumerate}
  \item 先看第 1 节的总览，明确作业暴露出的薄弱点。
  \item 再看第 2 节的“模板迁移法”，尤其是 BFS、Dijkstra、MST 和树递归。
  \item 最后逐题看第 3 节，训练“看到题面如何把 md 里的模板改成能 AC 的代码”。
\end{enumerate}

\begin{focusbox}
这份资料的中心结论是：你不是不会模板，而是需要把模板拆成四个可替换部件：\textbf{状态、转移、判重、答案口径}。只要这四件事判断对，大多数题都能从已有模板自然改出来。
\end{focusbox}

\section{作业体现的知识点与个人薄弱点}

\subsection{总体知识地图}
本次 PDF 作业一共覆盖 21 道题，可以分为两大块：树结构题与图搜索题。树题偏向“输入形式如何转成遍历/递归”，图题偏向“普通 BFS 什么时候需要升级或改状态”。

\begin{longtable}{L{2.9cm}L{4.2cm}L{7.2cm}}
\toprule
\textbf{模块} & \textbf{涉及题目} & \textbf{真正考点} \\
\midrule
二叉树遍历与重建 & 22、23、49、50、53 & 前序/中序/后序之间的递归关系；扩展先序和括号嵌套输入的解析；空树出口如何设定。\\
BST 与完全二叉树编号 & 24、25、54 & BST 插入与验证；重复值处理；完全二叉树编号不建树，用区间计算。\\
堆与 Huffman & 26、28 & \texttt{heapq} 的小顶堆性质；每组初始化；Huffman 的答案是每次合并代价之和。\\
BFS 与状态 BFS & 单词序列、仙岛求药、变换的迷宫、打怪救公主 & 是否普通无权；状态是否要包含时间、资源、剩余血量；判重是否还能用布尔数组。\\
Dijkstra 与路径计数 & 拯救行动、A Walk Through the Forest & 边权不同则不能普通 BFS；最短路结果可以作为后续 DP 的评价函数。\\
Floyd 与路径恢复 & 兔子与樱花 & 多源多查询；不仅要最短距离，还要恢复完整路径并按格式输出。\\
MST 与拓扑 & 拓扑排序、丛林中的路、连接所有点的最小费用 & DAG 入度维护；编号小优先用最小堆；显式边用 Kruskal，完全图用 Prim。\\
\bottomrule
\end{longtable}

\subsection{最明显的个人薄弱点}

\subsubsection{薄弱点一：模板能读懂，但题目一变就不知道状态该怎么改}
\texttt{graph.md} 里的 BFS 模板是标准的：
\begin{lstlisting}[style=py]
q = deque([(sx, sy, dist)])
visited[sx][sy] = True
while q:
    x, y, d = q.popleft()
    for dx, dy in dirs:
        ...
\end{lstlisting}
但是作业中真正容易错的地方是：状态不一定只有 \texttt{(x, y)}。例如：
\begin{itemize}
  \item 变换的迷宫：同一个格子在不同时间模数下可走性不同，所以状态是 \texttt{(x, y, t \% K)}。
  \item 打怪救公主：同一个格子如果以更多血量到达，后续更有希望，所以不能简单 visited，要用 \texttt{best[x][y]}。
  \item 拯救行动：进入不同格子的代价不同，普通队列按层扩展失效，要换成优先队列。
\end{itemize}

\subsubsection{薄弱点二：容易把“最短路、路线数、最小生成树”混在一起}
这三类题长得都像图题，但目标函数完全不同：
\begin{itemize}
  \item 最短路：问从 A 到 B 或从某点到所有点的最小距离。
  \item 路线数：问满足某种合法条件的路径有多少条，通常要 DFS/DP 计数。
  \item 最小生成树：问把所有点连通且总代价最小，不关心从某点到某点的最短距离。
\end{itemize}
作业中的典型对比：
\begin{itemize}
  \item 拯救行动是最短路，答案是最短时间。
  \item A Walk Through the Forest 先算最短路，但最终答案是合法路线数量。
  \item 丛林中的路和连接所有点的最小费用都是 MST，不是 Dijkstra。
\end{itemize}

\subsubsection{薄弱点三：树题容易过度建模}
\texttt{tree.md} 中很多代码使用节点类，这有助于理解。但作业里很多题并不要求真的建树：
\begin{itemize}
  \item 22 前序 + 中序输出后序：可以直接递归字符串返回后序。
  \item 24 完全二叉树编号：完全不用树节点，只需每层编号区间。
  \item 28 Huffman WPL：只需要堆和累计代价，不必真的构造编码树。
\end{itemize}
考试中要先问自己：题目要输出结构，还是只要一个序列/数字？如果只要序列或数字，常常可以用更轻的实现。

\subsubsection{薄弱点四：输入输出口径容易被忽略}
很多错题不是算法本身错，而是口径错：
\begin{itemize}
  \item 起点距离是 0 还是 1：问“时间”通常从 0；问“经过格子数”可能从 1。
  \item 多组输入如何结束：EOF、单独 0、0 0、先给 T，四种都出现过。
  \item 图是否无向：丛林中的路、兔子与樱花、A Walk Through the Forest 都要双向加边。
  \item 输出格式：拓扑排序是否带 \texttt{v} 前缀，兔子与樱花是否要 \texttt{A->(d)->B}。
\end{itemize}

\section{从 md 模板到实例代码的迁移方法}

\subsection{BFS 模板迁移：先问四个问题}
\begin{methodbox}
把任何 BFS 题都拆成四个问题：
\begin{enumerate}
  \item 状态是什么？
  \item 从一个状态能转移到哪些状态？
  \item 什么情况下两个状态算“已经访问过”？
  \item 答案是第一次到达、最短距离、最大资源，还是路径数量？
\end{enumerate}
\end{methodbox}

\subsubsection{普通 BFS}
适用条件：每条边代价相同，第一次到达某状态就是最优。
\begin{lstlisting}[style=py]
from collections import deque

q = deque([(sx, sy, 0)])
visited = [[False] * C for _ in range(R)]
visited[sx][sy] = True

while q:
    x, y, d = q.popleft()
    if (x, y) == (tx, ty):
        print(d)
        break
    for dx, dy in dirs:
        nx, ny = x + dx, y + dy
        if not in_bound(nx, ny):
            continue
        if wall(nx, ny) or visited[nx][ny]:
            continue
        visited[nx][ny] = True
        q.append((nx, ny, d + 1))
\end{lstlisting}

作业对应：仙岛求药、单词序列的 BFS 部分。注意单词序列里状态不是坐标，而是单词；仙岛求药里答案可能是经过格子数，起点距离可能设为 1。

\subsubsection{状态 BFS}
当同一个位置在不同附加状态下有不同意义时，visited 必须升维。
\begin{lstlisting}[style=py]
visited = [[[False] * K for _ in range(C)] for _ in range(R)]
q = deque([(sx, sy, 0)])
visited[sx][sy][0] = True

while q:
    x, y, t = q.popleft()
    nt = t + 1
    key = nt % K
    ...
    if not visited[nx][ny][key]:
        visited[nx][ny][key] = True
        q.append((nx, ny, nt))
\end{lstlisting}

作业对应：变换的迷宫。它的关键不是“迷宫 BFS”，而是“墙随时间变化，所以同一格不同时间模数是不同状态”。

\subsubsection{最优值 BFS}
当题目目标不是最短步数，而是“剩余资源最大”时，布尔 visited 会误杀好状态。
\begin{lstlisting}[style=py]
best = [[-1] * C for _ in range(R)]
best[sx][sy] = initial_hp
q = deque([(sx, sy)])

while q:
    x, y = q.popleft()
    cur = best[x][y]
    for dx, dy in dirs:
        nx, ny = x + dx, y + dy
        nhp = cur - cost(nx, ny)
        if nhp <= 0:
            continue
        if nhp > best[nx][ny]:
            best[nx][ny] = nhp
            q.append((nx, ny))
\end{lstlisting}

作业对应：打怪救公主。你要把“是否访问过”改成“是否以更好的血量访问过”。

\subsubsection{Dijkstra：当 BFS 不再成立}
只要进入不同邻居的代价不同，普通 BFS 的“按层最短”就失效。
\begin{lstlisting}[style=py]
import heapq

dist = [[float('inf')] * C for _ in range(R)]
dist[sx][sy] = 0
pq = [(0, sx, sy)]

while pq:
    d, x, y = heapq.heappop(pq)
    if d != dist[x][y]:
        continue
    if (x, y) == (tx, ty):
        print(d)
        break
    for dx, dy in dirs:
        nx, ny = x + dx, y + dy
        w = cost(nx, ny)
        nd = d + w
        if nd < dist[nx][ny]:
            dist[nx][ny] = nd
            heapq.heappush(pq, (nd, nx, ny))
\end{lstlisting}

作业对应：拯救行动。普通格代价 1，守卫格代价 2，所以要用优先队列。

\subsection{Dijkstra 模板迁移：不只是“从起点到终点”}
在 \texttt{graph.md} 里，Dijkstra 是从一个起点到所有点。作业中有两个重要变形。

\subsubsection{网格 Dijkstra}
拯救行动把网格看作图：
\begin{itemize}
  \item 每个格子是一个顶点。
  \item 上下左右相邻是边。
  \item 进入 \texttt{x} 的代价是 2，进入普通格/终点的代价是 1。
\end{itemize}
所以 Dijkstra 的图不一定要显式存邻接表，可以在扩展当前格子时动态枚举四个方向。

\subsubsection{反向 Dijkstra + DP}
A Walk Through the Forest 的关键是“每一步更接近家”。“更接近家”不是边权更小，而是到家的最短距离更小。于是：
\begin{enumerate}
  \item 从家 2 出发跑 Dijkstra，得到每个点到家的最短距离 \texttt{dist[u]}。
  \item 从办公室 1 开始 DFS。
  \item 只能走向满足 \texttt{dist[v] < dist[u]} 的邻居。
  \item 用 memo 记忆化每个点到家的合法路线数。
\end{enumerate}

\begin{lstlisting}[style=py]
def count(u):
    if u == 2:
        return 1
    if memo[u] != -1:
        return memo[u]
    ans = 0
    for v, w in graph[u]:
        if dist[v] < dist[u]:
            ans += count(v)
    memo[u] = ans
    return ans
\end{lstlisting}

\begin{warnbox}
这题不是求最短路径条数。只要每一步到家的最短距离严格下降，路径就合法，即使这条路径本身不是从 1 到 2 的最短路径。
\end{warnbox}

\subsection{Floyd 模板迁移：距离矩阵之外还要路径矩阵}
兔子与樱花是典型的“点数小 + 多次查询 + 要输出路径”的题。只用 Dijkstra 也可以，但每次查询都跑会麻烦；Floyd 一次预处理全部点对更自然。

\begin{lstlisting}[style=py]
dist = [[INF] * n for _ in range(n)]
nxt = [[-1] * n for _ in range(n)]
for i in range(n):
    dist[i][i] = 0
    nxt[i][i] = i

for u, v, w in edges:
    if w < dist[u][v]:
        dist[u][v] = dist[v][u] = w
        nxt[u][v] = v
        nxt[v][u] = u

for k in range(n):
    for i in range(n):
        for j in range(n):
            if dist[i][k] + dist[k][j] < dist[i][j]:
                dist[i][j] = dist[i][k] + dist[k][j]
                nxt[i][j] = nxt[i][k]
\end{lstlisting}

路径恢复：
\begin{lstlisting}[style=py]
def get_path(u, v):
    if nxt[u][v] == -1:
        return []
    path = [u]
    while u != v:
        u = nxt[u][v]
        path.append(u)
    return path
\end{lstlisting}

\subsection{MST 模板迁移：Prim 和 Kruskal 的选择}
\begin{longtable}{L{3cm}L{4cm}L{6.5cm}}
\toprule
\textbf{场景} & \textbf{优先算法} & \textbf{原因} \\
\midrule
边已经显式给出，例如 A B 12 & Kruskal & 直接收集边、排序、并查集判断是否成环。\\
任意两点都能连边，边权由坐标计算 & Prim & 不必生成所有边；维护每个未选点到已选集合的最小连接代价。\\
题目强调“连通且总代价最小” & MST & 不要误用 Dijkstra；Dijkstra 是从某点出发的最短路。\\
\bottomrule
\end{longtable}

Kruskal 核心：
\begin{lstlisting}[style=py]
edges.sort()
parent = list(range(n))

def find(x):
    if parent[x] != x:
        parent[x] = find(parent[x])
    return parent[x]

def union(a, b):
    ra, rb = find(a), find(b)
    if ra == rb:
        return False
    parent[rb] = ra
    return True

ans = 0
cnt = 0
for w, u, v in edges:
    if union(u, v):
        ans += w
        cnt += 1
        if cnt == n - 1:
            break
\end{lstlisting}

Prim 核心：
\begin{lstlisting}[style=py]
used = [False] * n
minDist = [float('inf')] * n
minDist[0] = 0
ans = 0

for _ in range(n):
    u = -1
    for i in range(n):
        if not used[i] and (u == -1 or minDist[i] < minDist[u]):
            u = i
    used[u] = True
    ans += minDist[u]

    for v in range(n):
        if not used[v]:
            w = weight(u, v)
            if w < minDist[v]:
                minDist[v] = w
\end{lstlisting}

\subsection{树模板迁移：输入形式决定递归方式}

\begin{longtable}{L{3.5cm}L{4.2cm}L{6cm}}
\toprule
\textbf{输入形式} & \textbf{根的位置} & \textbf{切分方法} \\
\midrule
前序 + 中序 & 前序第一个 & 在中序找根；左子树长度决定前序切片。\\
中序 + 后序 & 后序最后一个 & 在中序找根；左子树长度决定后序切片。\\
扩展先序 & 当前字符 & 当前字符是 \texttt{.} 则为空；否则先建左再建右。\\
括号嵌套 & 子串第一个字符 & 括号内部按最外层逗号切左右子树。\\
编号左右孩子 & 根通常由题面给出 & 用数组保存左右孩子，递归传编号。\\
\bottomrule
\end{longtable}

\section{逐题详细复盘：作业与 md 的拓展关系}

\subsection{22 根据二叉树前中序序列建树}
\begin{focusbox}
md 来源：\texttt{tree.md} 中“重建二叉树”的前序 + 中序原则。作业拓展：多组输入到 EOF，输出后序。
\end{focusbox}

模板迁移：
\begin{enumerate}
  \item md 中说明 \texttt{pre[0]} 是根。
  \item 在中序中找到根的位置 \texttt{k}。
  \item 中序左半边长度就是左子树节点数。
  \item 前序切片不能随便切，必须按左子树长度切。
  \item 输出后序，所以递归结果是 \texttt{left + right + root}。
\end{enumerate}

\begin{warnbox}
本题最容易错在前序切片：左子树前序不是 \texttt{pre[1:k]}，而是 \texttt{pre[1:1+len(left\_in)]}。因为 \texttt{k} 是根在中序中的位置，不是前序中的位置。
\end{warnbox}

可背思路：
\begin{lstlisting}[style=py]
def post(pre, ino):
    if not pre:
        return ''
    root = pre[0]
    k = ino.index(root)
    left_in = ino[:k]
    right_in = ino[k+1:]
    left_pre = pre[1:1+len(left_in)]
    right_pre = pre[1+len(left_in):]
    return post(left_pre, left_in) + post(right_pre, right_in) + root
\end{lstlisting}

\subsection{23 括号嵌套二叉树}
md 中主要是普通树和二叉树遍历，本题补充的是“字符串解析”。如果直接 \texttt{split(',')}，遇到 \texttt{A(B(*,C),D(E))} 会把内部逗号也切掉。

迁移方法：
\begin{enumerate}
  \item 当前子串为空或 \texttt{*}，返回空树。
  \item 当前子串第一个字符是根。
  \item 如果长度为 1，就是叶子。
  \item 去掉外层根和括号，得到内部字符串。
  \item 扫描内部字符串，维护括号深度；只在深度为 0 的逗号处分割。
\end{enumerate}

核心代码：
\begin{lstlisting}[style=py]
def split_children(s):
    depth = 0
    for i, ch in enumerate(s):
        if ch == '(':
            depth += 1
        elif ch == ')':
            depth -= 1
        elif ch == ',' and depth == 0:
            return s[:i], s[i+1:]
    return s, ''
\end{lstlisting}

这题是“树 + 解析”的组合题。复习时不要只背遍历，还要练习如何把题面给的字符串表示转成左右子树。

\subsection{24 二叉树}
\begin{focusbox}
md 来源：完全二叉树的编号性质：节点 \texttt{p} 的左孩子是 \texttt{2p}，右孩子是 \texttt{2p+1}。作业拓展：求编号 \texttt{m} 所在子树实际存在的节点数。
\end{focusbox}

不要建树。因为 \texttt{n} 可很大，建树既慢又浪费。以 \texttt{m} 为根的子树每一层编号都是连续区间：
\begin{lstlisting}[style=plain]
第0层: [m, m]
第1层: [2m, 2m+1]
第2层: [4m, 4m+3]
下一层: left *= 2, right = right * 2 + 1
\end{lstlisting}

每层贡献：
\begin{lstlisting}[style=py]
ans += min(right, n) - left + 1
\end{lstlisting}
前提是 \texttt{left <= n}，否则这一层已经没有实际存在的节点。

\subsection{25 二叉搜索树的层次遍历}
md 中有 BST 插入、树节点、层序 BFS。作业把它们拼起来，额外强调重复值只插入一次。

迁移方法：
\begin{enumerate}
  \item 输入一串数字。
  \item 按顺序插入 BST。
  \item 插入时小于走左，大于走右，等于直接忽略。
  \item 建好树后用 \texttt{deque} 层序遍历。
  \item 输出时用 \texttt{' '.join(ans)}，避免末尾空格。
\end{enumerate}

高频错误：
\begin{itemize}
  \item 把重复值插入左边或右边，导致层序结果多数字。
  \item 用 \texttt{list.pop(0)} 虽然能跑小数据，但 \texttt{deque.popleft()} 更规范。
  \item 空输入时访问空根，不过本题通常保证有输入。
\end{itemize}

\subsection{26 实现堆结构}
md 中既有手写二叉堆，也有 Python \texttt{heapq}。作业更适合直接用 \texttt{heapq}。

关键迁移：
\begin{itemize}
  \item 题目要求“输出并删除最小元素”，而 \texttt{heapq} 默认就是最小堆。
  \item 每一组测试数据都要重新 \texttt{heap=[]}。
  \item \texttt{type=1} 后面有一个数，\texttt{type=2} 没有额外参数。
\end{itemize}

API 检查：
\begin{lstlisting}[style=py]
import heapq
heapq.heappush(heap, x)
x = heapq.heappop(heap)
heapq.heapify(arr)
\end{lstlisting}

\subsection{27 树的转换}
md 中讲“左儿子右兄弟”转换：第一个孩子变左儿子，其余孩子串成右兄弟。作业问高度，不要求输出树。

理解重点：
\begin{itemize}
  \item 原树高度：只看父子关系向下走了多少层。
  \item 转换后二叉树高度：左孩子边和右兄弟边都算二叉树的边。
  \item 所以兄弟多时，转换后二叉树高度可能显著增加。
\end{itemize}

本题真正考的是“能否只维护答案，不真的构造结构”。看到只求高度，可以用栈模拟 DFS 过程，把当前节点对应的二叉深度存在栈里。

\subsection{28 Huffman 编码树}
md 中讲了完整 Huffman 树构造和编码，作业只要求最小 WPL。

最小 WPL 的核心解释：
\begin{itemize}
  \item 每次合并两个最小权值 \texttt{a} 和 \texttt{b}。
  \item 这次合并会让这两个子树里的所有叶子深度加 1。
  \item 增加的 WPL 正好是 \texttt{a+b}。
  \item 所以答案累计所有合并值。
\end{itemize}

不要把答案写成最后堆中根权值。最后根权值只是所有叶子权值总和，不是 WPL。

\subsection{49 二叉树的深度}
md 中的深度递归是 \texttt{max(left,right)+1}。作业输入不是节点对象，而是编号描述：
\begin{itemize}
  \item 第 \texttt{i} 行给节点 \texttt{i} 的左右孩子。
  \item \texttt{-1} 表示空。
  \item 根节点明确是 1。
\end{itemize}

所以把 \texttt{depth(root)} 改成 \texttt{depth(i)}：
\begin{lstlisting}[style=py]
def depth(i):
    if i == -1:
        return 0
    return max(depth(left[i]), depth(right[i])) + 1
\end{lstlisting}

\subsection{50 扩展二叉树}
扩展二叉树输入是先序序列，\texttt{.} 表示空节点。它和普通先序的区别是：普通先序无法唯一确定树，但带空标记的先序可以。

迁移步骤：
\begin{enumerate}
  \item 用全局下标 \texttt{pos} 指向当前读到的字符。
  \item 读一个字符后 \texttt{pos += 1}。
  \item 如果字符是 \texttt{.}，返回空。
  \item 否则创建节点，并递归构建左子树、右子树。
  \item 最后输出中序和后序，遍历时忽略空节点。
\end{enumerate}

\subsection{53 根据二叉树中后序序列建树}
它和第 22 题是镜像关系。前中题根在前序开头；中后题根在后序末尾。

可背思路：
\begin{lstlisting}[style=py]
def preorder(ino, post):
    if not ino:
        return ''
    root = post[-1]
    k = ino.index(root)
    left_in = ino[:k]
    right_in = ino[k+1:]
    left_post = post[:len(left_in)]
    right_post = post[len(left_in):-1]
    return root + preorder(left_in, left_post) + preorder(right_in, right_post)
\end{lstlisting}

高频错误：后序切右子树时不要包含最后的根。

\subsection{54 验证二叉搜索树}
md 中讲 BST 插入和查找，作业问“给定树是否是 BST”。判断 BST 合法不能只看父子。

反例：
\begin{lstlisting}[style=plain]
    10
   /  \
  5    15
      /
     6
\end{lstlisting}
节点 6 小于 15，看起来是 15 的合法左孩子；但它在 10 的右子树里，必须大于 10，所以整棵树不是 BST。

正确方法是上下界递归：
\begin{lstlisting}[style=py]
def valid(node, low, high):
    if node is None:
        return True
    if not (low < node.val < high):
        return False
    return valid(node.left, low, node.val) and valid(node.right, node.val, high)
\end{lstlisting}

\subsection{1 单词序列}
\texttt{graph.md} 的词梯问题就是本题基础。迁移时注意三件事：
\begin{enumerate}
  \item 点是单词，不是数字编号。
  \item 两个单词只差一个字母时有边。
  \item 如果输出的是变换序列长度，起点长度通常是 1，不是 0。
\end{enumerate}

建图可以两种方式：
\begin{itemize}
  \item 暴力比较每两个单词，简单但可能慢。
  \item 桶方法：把每个位置替换成 \texttt{\_}，同桶单词互相可达。
\end{itemize}

\subsection{2 仙岛求药}
这是普通网格 BFS，但要特别注意答案口径。

迁移检查：
\begin{itemize}
  \item 地图中的起点和终点字符是什么？
  \item 墙是什么字符？
  \item 输出的是移动时间，还是经过格子数？
  \item 不可达输出什么？
\end{itemize}

如果题目说“经过的方格数包括起点”，起点距离设为 1。如果题目说“花费时间”，起点时间设为 0。

\subsection{3 变换的迷宫}
这是状态 BFS 的重点题。普通 \texttt{visited[x][y]} 会错，因为同一个格子在不同时间可能可走性不同。

迁移方法：
\begin{enumerate}
  \item 状态写成 \texttt{(x, y, t)}。
  \item 判重写成 \texttt{visited[x][y][t \% K]}。
  \item 每走一步，时间 \texttt{t+1}。
  \item 判断下一格是否可走时，要用下一时刻的时间模数。
\end{enumerate}

记忆句：只要题目出现周期、时间变化、白天黑夜、机关开关，就先想 \texttt{time \% period}。

\subsection{1 拓扑排序}
\texttt{graph.md} 中拓扑排序用普通队列即可。但作业里常见额外要求：同等条件下编号小的点先输出。

迁移：
\begin{itemize}
  \item 如果只要求任意拓扑序：用 \texttt{deque}。
  \item 如果要求编号小优先：用 \texttt{heapq} 存入度为 0 的点。
  \item 如果最后输出点数不足 \texttt{n}：说明有环。
\end{itemize}

核心逻辑：
\begin{lstlisting}[style=py]
for v in graph[u]:
    indeg[v] -= 1
    if indeg[v] == 0:
        heapq.heappush(heap, v)
\end{lstlisting}

\subsection{2 丛林中的路}
这是 Kruskal MST。md 里讲的并查集在这里完整使用。

作业特殊点：
\begin{itemize}
  \item 顶点是大写字母，需要转成编号：\texttt{ord(ch)-ord('A')}。
  \item 输入每行可能包含多个邻接边，要循环读。
  \item 是无向图，但 Kruskal 只需要边集合，不必存双向邻接表。
\end{itemize}

\subsection{3 A Walk Through the Forest}
本题是组合题，非常适合作为考前重点。

题面含义：
\begin{itemize}
  \item 家是 2，办公室是 1。
  \item Jimmy 每一步必须更接近家。
  \item “更接近”指到家 2 的最短距离更小。
\end{itemize}

算法组合：
\begin{enumerate}
  \item 建无向带权图。
  \item 从 2 跑 Dijkstra，得到 \texttt{dist[u]}。
  \item 从 1 做 DFS 计数，只能走向 \texttt{dist[v] < dist[u]} 的点。
  \item 用 \texttt{memo[u]} 避免重复计算。
\end{enumerate}

\begin{warnbox}
这题最容易误解成“求从 1 到 2 的最短路径条数”。实际上，只要每一步越来越接近家，整条路径就合法，它可能不是最短路径。
\end{warnbox}

\subsection{4 打怪救公主}
本题问“到达公主时最大剩余血量”，所以不是普通最短路题。

模板改法：
\begin{itemize}
  \item 普通 BFS 的 \texttt{visited[x][y]} 改成 \texttt{best[x][y]}。
  \item \texttt{best[x][y]} 表示目前已知到达该格子的最大剩余血量。
  \item 如果新路径剩余血量更大，即使这个格子以前来过，也要重新入队。
\end{itemize}

判断：
\begin{itemize}
  \item 进入数字格扣血。
  \item 起点、终点、普通路不扣血。
  \item 血量扣到 0 或以下不能继续。
\end{itemize}

\subsection{5 兔子与樱花}
这是 Floyd + 路径恢复 + 字符串映射。

完整迁移：
\begin{enumerate}
  \item 读地点名，建立 \texttt{id\_of[name]}。
  \item 初始化 \texttt{dist} 和 \texttt{nxt}。
  \item 读无向边，更新两边。
  \item Floyd 三重循环更新最短路。
  \item 查询时恢复点序列。
  \item 按 \texttt{A->(d)->B} 输出每一段。
\end{enumerate}

注意：输出每段的 \texttt{d} 是路径上相邻两点的边长，不是从起点到该点的累计距离。

\subsection{56 拯救行动}
这是网格 Dijkstra。它和打怪救公主看起来都像“救人迷宫”，但目标不同：
\begin{itemize}
  \item 打怪救公主：最大剩余血量。
  \item 拯救行动：最短时间。
\end{itemize}

所以拯救行动要用 \texttt{dist[x][y]}，表示到达该格子的最短时间。

代价：
\begin{itemize}
  \item 进入 \texttt{@}、\texttt{r}、\texttt{a}：代价 1。
  \item 进入 \texttt{x}：代价 2。
  \item 进入 \texttt{\#}：不允许。
\end{itemize}

\subsection{57 连接所有点的最小费用}
这是完全图上的 MST。

题面信号：
\begin{itemize}
  \item “连接所有点”
  \item “总费用最小”
  \item “任意两点之间有且仅有一条简单路径”
\end{itemize}
这三个信号都指向最小生成树。

为什么用 Prim：
\begin{itemize}
  \item 任意两点都有边，是完全图。
  \item 如果建所有边再排序，需要 $O(n^2)$ 条边。
  \item Prim 可以边算曼哈顿距离边更新 \texttt{minDist}。
\end{itemize}

\section{考试时的算法选择决策表}

\begin{longtable}{L{4.2cm}L{4cm}L{6cm}}
\toprule
\textbf{题面信号} & \textbf{选择} & \textbf{理由与改法} \\
\midrule
每一步代价相同，问最少步数 & BFS & 队列按层扩展，第一次到达终点最优。\\
地图/规则随时间周期变化 & 状态 BFS & 状态加入 \texttt{t\%K}，visited 升维。\\
同一位置带不同资源 & 最优值 BFS 或状态 BFS & 资源影响未来，布尔 visited 不够。\\
进入不同点代价不同 & Dijkstra & 普通 BFS 按步数，不按总代价。\\
多次询问任意两点最短路 & Floyd & 点数小，用三重循环一次预处理。\\
要求输出最短路径本身 & BFS parent / Floyd nxt / Dijkstra pred & 只存距离不够，要保存路径恢复信息。\\
所有点连通且总代价最小 & MST & 与从某点出发最短路无关。\\
显式给边的 MST & Kruskal & 排序边 + 并查集。\\
完全图坐标 MST & Prim & 不建全边，动态算边权。\\
有向依赖顺序 & 拓扑排序 & 入度为 0 的点逐步输出。\\
编号小优先的拓扑序 & 堆优化拓扑 & 入度 0 的集合用最小堆。\\
前序 + 中序 & 树递归 & 根在前序开头，中序切左右。\\
中序 + 后序 & 树递归 & 根在后序末尾，中序切左右。\\
Huffman / 最优二叉树 / WPL & 最小堆 & 每次合并两个最小，累计合并值。\\
\bottomrule
\end{longtable}

\section{代码落地前的详细检查清单}

\begin{checkbox}
\begin{enumerate}
  \item 输入规模是否允许建图/建树？如果 \texttt{n} 很大，避免显式建树或建完全图边集。
  \item 输入结束条件是什么？EOF、\texttt{0}、\texttt{0 0}、还是第一行 \texttt{T}。
  \item 起点答案初始化是 0 还是 1？看题目问“时间/步数”还是“经过格子数/序列长度”。
  \item 图是有向还是无向？无向图读边时两边都要加，Kruskal 边集则只存一份即可。
  \item 点编号从 0 还是 1 开始？若题目是 1 到 n，数组一般开 \texttt{n+1}。
  \item \texttt{visited} 是布尔吗？如果题目有时间、钥匙、血量、剩余次数，通常不是简单二维布尔。
  \item BFS 是否仍然成立？只要边权不是全相等，就考虑 Dijkstra。
  \item 是否要输出路径？要输出路径就要保存 \texttt{parent}/\texttt{nxt}/\texttt{pred}。
  \item 是否要求字典序或编号小优先？队列可能要换成堆。
  \item 输出格式有没有箭头、括号、前缀、空格、大小写固定要求？
\end{enumerate}
\end{checkbox}

\section{考前重点训练顺序}

\subsection{第一优先级：模板变形题}
\begin{enumerate}
  \item 变换的迷宫：训练状态 BFS 和 \texttt{visited[x][y][t\%K]}。
  \item 拯救行动：训练普通 BFS 升级 Dijkstra。
  \item 打怪救公主：训练布尔 visited 改最优值数组。
  \item A Walk Through the Forest：训练 Dijkstra + 记忆化 DFS 计数。
  \item 兔子与樱花：训练 Floyd + 路径恢复 + 字符串映射。
\end{enumerate}

\subsection{第二优先级：树输入解析题}
\begin{enumerate}
  \item 22 与 53 成对复习，区分根在前序开头还是后序末尾。
  \item 23 单独复习括号嵌套，重点是最外层逗号。
  \item 50 单独复习扩展先序，重点是 \texttt{.} 的空节点出口。
  \item 54 单独复习 BST 验证，重点是上下界。
\end{enumerate}

\subsection{第三优先级：MST 与堆}
\begin{enumerate}
  \item 丛林中的路：显式边 Kruskal。
  \item 连接所有点的最小费用：完全图 Prim。
  \item 实现堆结构：\texttt{heapq} 操作流。
  \item Huffman 编码树：堆合并与 WPL。
\end{enumerate}

\section{最后的压缩记忆}

\begin{focusbox}
\begin{itemize}
  \item BFS 不是万能的；它只保证“边权相同”时第一次到达最优。
  \item 状态不止位置；时间、血量、钥匙、剩余次数都可能是状态的一部分。
  \item 最短路不是 MST；“连通所有点总代价最小”才是 MST。
  \item 树题先看输入形式，再决定递归根在哪里。
  \item 题目只问数字或序列时，不必执着于完整建模。
  \item 输出格式和输入结束条件是最容易白丢分的地方。
\end{itemize}
\end{focusbox}

\end{document}
